\definecolor{dsdlcodebg}{RGB}{246,248,250}
\definecolor{dsdlcodeborder}{RGB}{208,215,222}
\newsavebox{\dsdlcodebox}
\newenvironment{dsdlcodeblock}
{%
  \par\noindent
  \begin{latin}
  \begin{lrbox}{\dsdlcodebox}
  \begin{minipage}{0.94\linewidth}
  \small
}
{%
  \end{minipage}
  \end{lrbox}
  \setlength{\fboxsep}{6pt}
  \fcolorbox{dsdlcodeborder}{dsdlcodebg}{\usebox{\dsdlcodebox}}
  \end{latin}
  \par
}

‫% -------------------------------------------------------
‫%  Abstract
‫% -------------------------------------------------------
‫
‫\بدون‌تورفتگی \مهم{چکیده}: در این آزمایش دو نوع مقایسه‌کننده با زبان توصیف سخت‌افزار
‫\لر{Verilog} و با استفاده از توصیف جریان داده طراحی و شبیه‌سازی شد. در بخش نخست، ابتدا
‫یک مقایسه‌کننده‌ی یک‌بیتی با سه خروجی بزرگ‌تر، مساوی و کوچک‌تر ساخته شد و سپس با اتصال
‫چهار نمونه از این پیمانه، یک مقایسه‌کننده‌ی چهاربیتی ترکیبی و سلسله‌مراتبی به دست آمد.
‫در بخش دوم، یک مقایسه‌کننده‌ی دنباله‌ای ترتیبی طراحی شد که پس از بازنشانی، بیت‌های دو عدد
‫را از پرارزش‌ترین بیت به کم‌ارزش‌ترین بیت دریافت می‌کند و نتیجه‌ی مقایسه‌ی پیشوند دریافت‌شده
‫را پس از هر نبض ساعت در خروجی نشان می‌دهد. مطابق محدودیت صورت آزمایش، پیمانه‌های طراحی
‫فقط از دستور \کد{assign} برای توصیف منطق استفاده می‌کنند و در مقایسه‌کننده‌ی دنباله‌ای نیز
‫ذخیره‌سازی حالت به کمک بازخورد پیوسته‌ی ساختار ارباب--برده انجام شده است. برای ارزیابی، هر
‫دو مدار با نرم‌افزار \لر{Icarus Verilog 12.0} شبیه‌سازی شدند. تمام ۲۵۶ حالت ممکن ورودی
‫چهاربیتی در مدار موازی و تمام ۲۵۶ زوج ورودی در مدار دنباله‌ای، در هر چهار مرحله‌ی دریافت بیت،
‫بدون خطا آزموده شدند.
‫
‫
‫%\پرش‌بلند
‫\بدون‌تورفتگی \مهم{واژه‌های کلیدی}:
‫وریلاگ، مقایسه‌کننده، توصیف جریان داده، طراحی سلسله‌مراتبی، مدار ترکیبی، مدار ترتیبی،
‫آی‌کاروس وریلاگ
‫%\صفحه‌جدید
‫
‫
‫\section{مقدمه}
‫
‫مقایسه‌کننده‌ی دیجیتال مداری است که دو عدد دودویی را دریافت می‌کند و رابطه‌ی بین آن‌ها را
‫به صورت بزرگ‌تر بودن، مساوی بودن یا کوچک‌تر بودن مشخص می‌سازد. این مدار را می‌توان به شکل
‫موازی طراحی کرد تا تمام بیت‌ها را هم‌زمان بررسی کند، یا به صورت دنباله‌ای ساخت تا در هر نبض
‫ساعت فقط یک جفت بیت را دریافت کرده و نتیجه‌ی مقایسه‌ی انجام‌شده تا آن لحظه را نگه دارد.
‫
‫هدف این آزمایش، آشنایی عملی با توصیف جریان داده\پاورقی{Dataflow description}، طراحی
‫سلسله‌مراتبی\پاورقی{Hierarchical design}، تفاوت مدارهای ترکیبی و ترتیبی و همچنین نوشتن
‫آزمون خودکار برای یک مدار دیجیتال است. تمام معادلات منطقی مدارهای اصلی با دستور
‫\کد{assign} نوشته شده‌اند و صحت عملکرد آن‌ها با شبیه‌ساز \لر{Icarus Verilog} بررسی شده است.
‫
‫
‫\subsection{تعریف مسئله}
‫
‫این آزمایش شامل دو بخش است:
‫
‫\شروع{فقرات}
‫\فقره طراحی یک مقایسه‌کننده‌ی یک‌بیتی به روش جریان داده و استفاده از چهار نمونه‌ی آن برای
‫ساخت یک مقایسه‌کننده‌ی چهاربیتی ترکیبی و سلسله‌مراتبی.
‫\فقره طراحی یک مقایسه‌کننده‌ی دنباله‌ای ترتیبی در قالب تنها یک پیمانه، به گونه‌ای که پیمانه
‫در ابتدای کار بازنشانی شود، سپس بیت‌های ورودی را در هر نبض ساعت دریافت کند و نتیجه‌ی
‫مقایسه‌ی پیشوند دریافت‌شده را در خروجی نمایش دهد.
‫\فقره پیاده‌سازی تمام منطق پیمانه‌های اصلی فقط با توصیف جریان داده و بدون استفاده از
‫ساختارهای رفتاری مانند \کد{always}، \کد{initial} و \کد{case}.
‫\پایان{فقرات}
‫
‫
‫\subsection{اهداف آزمایش}
‫
‫اهداف اصلی انجام آزمایش عبارت‌اند از:
‫
‫\شروع{فقرات}
‫\فقره استخراج معادلات بولی یک مقایسه‌کننده‌ی یک‌بیتی؛
‫\فقره استفاده‌ی مجدد از یک پیمانه‌ی ساده در یک طراحی بزرگ‌تر؛
‫\فقره رعایت اولویت بیت پرارزش‌تر در مقایسه‌ی اعداد چندبیتی؛
‫\فقره طراحی حالت برای مقایسه‌ی دنباله‌ای و حفظ نتیجه‌ی قطعی‌شده؛
‫\فقره درک نقش بازنشانی و ساعت در مدار ترتیبی؛
‫\فقره آزمون جامع مدارها با تولید خودکار تمام ورودی‌های ممکن.
‫\پایان{فقرات}
‫
‫
‫\section{مفاهیم اولیه}
‫
‫
‫\subsection{توصیف جریان داده}
‫
‫در توصیف جریان داده، ارتباط میان ورودی و خروجی به وسیله‌ی عبارت‌های منطقی و دستور
‫\کد{assign} بیان می‌شود. سمت راست یک انتساب پیوسته هر زمان که یکی از ورودی‌های آن تغییر
‫کند دوباره ارزیابی می‌شود و نتیجه روی سیم خروجی قرار می‌گیرد. برای نمونه، عبارت زیر
‫ترکیب دروازه‌های \لر{AND} و \لر{NOT} را توصیف می‌کند:
‫
\begin{dsdlcodeblock}
\begin{verbatim}
10  assign gt =  a & ~b;
\end{verbatim}
\end{dsdlcodeblock}
‫
‫این روش برای بیان مدارهای ترکیبی بسیار مناسب است، زیرا خروجی مدار ترکیبی فقط به مقدار
‫فعلی ورودی‌ها وابسته است.
‫
‫
‫\subsection{مدار ترکیبی و مدار ترتیبی}
‫
‫در \مهم{مدار ترکیبی} هیچ حافظه‌ای وجود ندارد و خروجی فقط تابع ورودی جاری است. مقایسه‌کننده‌ی
‫چهاربیتی این آزمایش یک مدار ترکیبی است. در مقابل، \مهم{مدار ترتیبی} علاوه بر ورودی جاری،
‫به نتیجه یا حالت قبلی نیز وابسته است. مقایسه‌کننده‌ی دنباله‌ای باید به یاد داشته باشد که در
‫بیت‌های قبلی کدام عدد بزرگ‌تر بوده است؛ بنابراین یک مدار ترتیبی محسوب می‌شود.
‫
‫
‫\subsection{اصل اولویت بیت پرارزش‌تر}
‫
‫در مقایسه‌ی دو عدد دودویی، نخستین بیت متفاوت از سمت چپ نتیجه را تعیین می‌کند. برای مثال
‫در مقایسه‌ی $1010_2$ و $1001_2$، دو بیت نخست برابرند، اما در بیت بعدی مقدار عدد اول برابر
‫$1$ و مقدار عدد دوم برابر $0$ است. در نتیجه عدد اول بزرگ‌تر است و بیت پایانی دیگر اثری بر
‫نتیجه ندارد.
‫
‫
‫\section{طراحی مقایسه‌کننده‌ی چهاربیتی}
‫
‫
‫\subsection{مقایسه‌کننده‌ی یک‌بیتی}
‫
‫ورودی‌های پیمانه‌ی یک‌بیتی با $a$ و $b$ و خروجی‌های آن با $gt$، $eq$ و $lt$ نمایش داده
‫می‌شوند. جدول درستی مدار در جدول~\رجوع{جدول:درستی یک‌بیتی} آمده است.
‫
‫
‫\شروع{لوح}[ht]
‫\تنظیم‌ازوسط
‫\شرح{جدول درستی مقایسه‌کننده‌ی یک‌بیتی}
‫
‫\شروع{جدول}{|c|c|c|c|c|}
‫\خط‌پر
‫\سیاه $a$ & \سیاه $b$ & \سیاه $gt$ & \سیاه $eq$ & \سیاه $lt$ \\
‫\خط‌پر \خط‌پر
‫$0$ & $0$ & $0$ & $1$ & $0$ \\
‫$0$ & $1$ & $0$ & $0$ & $1$ \\
‫$1$ & $0$ & $1$ & $0$ & $0$ \\
‫$1$ & $1$ & $0$ & $1$ & $0$ \\
‫\خط‌پر
‫\پایان{جدول}
‫
‫\برچسب{جدول:درستی یک‌بیتی}
‫\پایان{لوح}
‫
‫
‫با توجه به جدول درستی، معادلات بولی سه خروجی به صورت زیر به دست می‌آیند:
‫
‫\شروع{equation}
‫gt = a \wedge \neg b
‫\پایان{equation}
‫
‫\شروع{equation}
‫eq = \neg(a \oplus b)
‫\پایان{equation}
‫
‫\شروع{equation}
‫lt = \neg a \wedge b
‫\پایان{equation}
‫
‫این معادلات در پرونده‌ی \کد{comparator\_1bit.v} مستقیماً در خطوط ۱۰ تا ۱۲ با سه دستور \کد{assign}
‫پیاده‌سازی شده‌اند:
‫
\begin{dsdlcodeblock}
\begin{verbatim}
10  assign gt =  a & ~b;
11  assign eq = ~(a ^ b);
12  assign lt = ~a &  b;
\end{verbatim}
\end{dsdlcodeblock}
‫
‫
‫\subsection{ساختار سلسله‌مراتبی مدار چهاربیتی}
‫
‫پیمانه‌ی \کد{comparator\_4bit.v} شامل چهار نمونه از مقایسه‌کننده‌ی یک‌بیتی است. نمونه‌ی
‫\کد{c3} پرارزش‌ترین بیت و نمونه‌ی \کد{c0} کم‌ارزش‌ترین بیت را مقایسه می‌کند. خطوط
‫۱۴ تا ۲۸ این پیمانه، چهار نمونه و منطق ترکیب خروجی‌ها را نشان می‌دهند. خروجی‌های هر نمونه در سه بردار میانی \کد{bit\_gt}، \کد{bit\_eq} و \کد{bit\_lt} ذخیره می‌شوند.
‫
\begin{dsdlcodeblock}
\begin{verbatim}
14  comparator_1bit c3(a[3], b[3], bit_gt[3], bit_eq[3], bit_lt[3]);
15  comparator_1bit c2(a[2], b[2], bit_gt[2], bit_eq[2], bit_lt[2]);
16  comparator_1bit c1(a[1], b[1], bit_gt[1], bit_eq[1], bit_lt[1]);
17  comparator_1bit c0(a[0], b[0], bit_gt[0], bit_eq[0], bit_lt[0]);
18
19  assign gt = bit_gt[3]
20            | (bit_eq[3] & bit_gt[2])
21            | (bit_eq[3] & bit_eq[2] & bit_gt[1])
22            | (bit_eq[3] & bit_eq[2] & bit_eq[1] & bit_gt[0]);
23  assign eq = &bit_eq;
24  assign lt = bit_lt[3]
25            | (bit_eq[3] & bit_lt[2])
26            | (bit_eq[3] & bit_eq[2] & bit_lt[1])
27            | (bit_eq[3] & bit_eq[2] & bit_eq[1] & bit_lt[0]);
28  endmodule
\end{verbatim}
\end{dsdlcodeblock}
‫
‫شکل~\رجوع{شکل:کد مقایسه‌کننده چهاربیتی} کد و ساختار سلسله‌مراتبی را نشان می‌دهد.
‫
‫\شروع{شکل}[ht]
‫\تنظیم‌ازوسط
‫\centerimg{1.png}{17cm}
‫\شرح{پیاده‌سازی جریان داده و ساختار سلسله‌مراتبی مقایسه‌کننده‌ی چهاربیتی}
‫\برچسب{شکل:کد مقایسه‌کننده چهاربیتی}
‫\پایان{شکل}
‫
‫
‫\subsection{استخراج خروجی‌های چهاربیتی}
‫
‫برای ساده‌سازی، نتیجه‌ی مقایسه‌ی هر بیت با $g_i$، $e_i$ و $l_i$ نمایش داده می‌شود. چون
‫بیت شماره‌ی ۳ پرارزش‌ترین بیت است، معادله‌ی بزرگ‌تر بودن به صورت زیر است:
‫
‫\شروع{equation}
‫GT = g_3 \vee (e_3 \wedge g_2)
‫\vee (e_3 \wedge e_2 \wedge g_1)
‫\vee (e_3 \wedge e_2 \wedge e_1 \wedge g_0)
‫\پایان{equation}
‫
‫به همین ترتیب، معادله‌ی کوچک‌تر بودن برابر است با:
‫
‫\شروع{equation}
‫LT = l_3 \vee (e_3 \wedge l_2)
‫\vee (e_3 \wedge e_2 \wedge l_1)
‫\vee (e_3 \wedge e_2 \wedge e_1 \wedge l_0)
‫\پایان{equation}
‫
‫دو عدد فقط زمانی برابرند که هر چهار جفت بیت با هم برابر باشند:
‫
‫\شروع{equation}
‫EQ = e_3 \wedge e_2 \wedge e_1 \wedge e_0
‫\پایان{equation}
‫
‫همان‌طور که در خط ۲۳ کد بالا دیده می‌شود، عملگر کاهش
‫\لر{AND} شرط برابری را به صورت کوتاه پیاده‌سازی می‌کند. بدین ترتیب خروجی تنها تابع ورودی
‫فعلی است و هیچ مسیر ذخیره‌سازی در مدار چهاربیتی وجود ندارد.
‫
‫
‫\section{طراحی مقایسه‌کننده‌ی دنباله‌ای}
‫
‫
‫\subsection{روش دریافت بیت‌ها}
‫
‫در این طراحی فرض شده است که بیت‌های دو عدد از پرارزش‌ترین بیت به کم‌ارزش‌ترین بیت وارد
‫می‌شوند. پیش از شروع مقایسه، ورودی \کد{reset} یک می‌شود تا حالت مدار به «برابر تا اینجا»
‫بازگردد. پس از غیرفعال شدن بازنشانی، در هر چرخه‌ی ساعت یک بیت روی ورودی $a$ و یک بیت روی
‫ورودی $b$ قرار می‌گیرد. نتیجه پس از لبه‌ی بالارونده‌ی ساعت در سه خروجی ظاهر می‌شود.
‫
‫
‫\subsection{کدگذاری حالت}
‫
‫برای نگه‌داری نتیجه‌ی مقایسه‌ی انجام‌شده، یک حالت دوبیتی در نظر گرفته شده است. جدول
‫حالت‌ها در جدول~\رجوع{جدول:حالت دنباله‌ای} آمده است.
‫
‫
‫\شروع{لوح}[ht]
‫\تنظیم‌ازوسط
‫\شرح{کدگذاری حالت مقایسه‌کننده‌ی دنباله‌ای}
‫
‫\شروع{جدول}{|c|c|c|c|c|}
‫\خط‌پر
‫\سیاه حالت & \سیاه مفهوم & \سیاه $gt$ & \سیاه $eq$ & \سیاه $lt$ \\
‫\خط‌پر \خط‌پر
‫\کد{00} & برابر تا این مرحله & $0$ & $1$ & $0$ \\
‫\کد{10} & عدد $a$ بزرگ‌تر است & $1$ & $0$ & $0$ \\
‫\کد{01} & عدد $a$ کوچک‌تر است & $0$ & $0$ & $1$ \\
‫\کد{11} & حالت استفاده‌نشده & -- & -- & -- \\
‫\خط‌پر
‫\پایان{جدول}
‫
‫\برچسب{جدول:حالت دنباله‌ای}
‫\پایان{لوح}
‫
‫
‫حالت \کد{00} برای بازنشانی انتخاب شده است. تا زمانی که دو بیت ورودی برابر باشند، حالت
‫تغییر نمی‌کند. اگر در حالت برابری $a=1$ و $b=0$ شود، مدار وارد حالت \کد{10} می‌شود. اگر
‫$a=0$ و $b=1$ شود، مدار وارد حالت \کد{01} می‌شود. پس از ورود به یکی از دو حالت بزرگ‌تر
‫یا کوچک‌تر، بیت‌های بعدی نتیجه را تغییر نمی‌دهند؛ زیرا آن‌ها کم‌ارزش‌ترند.
‫
‫
‫\subsection{منطق حالت بعدی}
‫
‫اگر $s_1$ نشان‌دهنده‌ی بزرگ‌تر بودن و $s_0$ نشان‌دهنده‌ی کوچک‌تر بودن باشد، منطق حالت
‫بعدی از روابط زیر پیروی می‌کند:
‫
‫\شروع{equation}
‫s_1^{+} = s_1 \vee (\neg s_1 \wedge \neg s_0 \wedge a \wedge \neg b)
‫\پایان{equation}
‫
‫\شروع{equation}
‫s_0^{+} = s_0 \vee (\neg s_1 \wedge \neg s_0 \wedge \neg a \wedge b)
‫\پایان{equation}
‫
‫وجود $s_1$ و $s_0$ در ابتدای عبارت‌های OR باعث می‌شود نتیجه‌ی قطعی‌شده حفظ شود. شرط
‫$\neg s_1 \wedge \neg s_0$ نیز اجازه می‌دهد ورودی جدید فقط در حالت برابری روی نتیجه اثر
‫بگذارد.
‫
‫
‫\subsection{ذخیره‌ی حالت با توصیف جریان داده}
‫
‫در توصیف متداول و قابل سنتز وریلاگ، حافظه‌ی لبه‌ای با یک بلوک \کد{always} نوشته می‌شود؛
‫اما صورت آزمایش استفاده از هر توصیف دیگری را ممنوع کرده است. به همین دلیل، حالت با دو
‫بخش ارباب و برده و بازخورد انتساب پیوسته مدل شده است. خطوط ۱۶ تا ۲۴ پرونده‌ی \کد{serial\_comparator.v} منطق حالت بعدی، ذخیره‌سازی و رمزگشایی خروجی را نشان می‌دهند:
‫
\begin{dsdlcodeblock}
\begin{verbatim}
16  assign next_state[1] = state[1] | (~state[1] & ~state[0] &  a & ~b);
17  assign next_state[0] = state[0] | (~state[1] & ~state[0] & ~a &  b);
18
19  assign master = reset ? 2'b00 : (~clk ? next_state : master);
20  assign state  = reset ? 2'b00 : ( clk ? master     : state);
21
22  assign gt =  state[1];
23  assign eq = ~state[1] & ~state[0];
24  assign lt =  state[0];
\end{verbatim}
\end{dsdlcodeblock}
‫
‫عملکرد این دو انتساب به ترتیب زیر است:
‫
‫\شروع{شمارش}
‫\فقره هنگامی که \کد{reset=1} است، هر دو بخش مقدار \کد{00} می‌گیرند.
‫\فقره هنگامی که \کد{clk=0} است، بخش ارباب مقدار حالت بعدی را دنبال می‌کند و بخش برده
‫مقدار قبلی را نگه می‌دارد.
‫\فقره هنگامی که \کد{clk=1} می‌شود، بخش ارباب مقدار خود را نگه می‌دارد و بخش برده مقدار
‫ذخیره‌شده در ارباب را دریافت می‌کند.
‫\فقره در نتیجه، خروجی قابل مشاهده در لبه‌ی بالارونده‌ی ساعت به‌روزرسانی می‌شود.
‫\پایان{شمارش}
‫
‫این روش محدودیت «فقط \کد{assign}» آزمایش را در شبیه‌ساز \لر{Icarus Verilog} برآورده
‫می‌کند. با این حال، بازخورد انتساب پیوسته روش متعارف برای سنتز مدار واقعی نیست و اگر محدودیت
‫آزمایش وجود نداشت، استفاده از یک فلیپ‌فلاپ لبه‌ای با توصیف ترتیبی انتخاب استانداردتری بود.
‫
‫شکل~\رجوع{شکل:کد دنباله‌ای} برای درج تصویر کامل پیمانه‌ی دنباله‌ای در نظر گرفته شده است.
‫
‫\شروع{شکل}[ht]
‫\تنظیم‌ازوسط
‫\centerimg{2.png}{15cm}
‫\شرح{پیاده‌سازی تک‌پیمانه‌ای و جریان داده‌ی مقایسه‌کننده‌ی دنباله‌ای}
‫\برچسب{شکل:کد دنباله‌ای}
‫\پایان{شکل}
‫
‫
‫\subsection{مثال مرحله‌به‌مرحله}
‫
‫برای نمایش عملکرد مدار دنباله‌ای، دو عدد $a=1010_2$ و $b=1001_2$ در نظر گرفته می‌شوند.
‫جدول~\رجوع{جدول:مثال دنباله‌ای} نتیجه‌ی مدار را پس از هر نبض ساعت نشان می‌دهد.
‫
‫
‫\شروع{لوح}[ht]
‫\تنظیم‌ازوسط
‫\شرح{مقایسه‌ی دنباله‌ای $1010_2$ و $1001_2$ از بیت پرارزش‌تر}
‫
‫\شروع{جدول}{|c|c|c|c|}
‫\خط‌پر
‫\سیاه مرحله & \سیاه بیت‌های ورودی $a,b$ & \سیاه پیشوند مقایسه‌شده & \سیاه نتیجه \\
‫\خط‌پر \خط‌پر
‫بازنشانی & -- & بدون بیت & مساوی \\
‫نبض ۱ & $1,1$ & $1$ و $1$ & مساوی \\
‫نبض ۲ & $0,0$ & $10$ و $10$ & مساوی \\
‫نبض ۳ & $1,0$ & $101$ و $100$ & بزرگ‌تر \\
‫نبض ۴ & $0,1$ & $1010$ و $1001$ & بزرگ‌تر \\
‫\خط‌پر
‫\پایان{جدول}
‫
‫\برچسب{جدول:مثال دنباله‌ای}
‫\پایان{لوح}
‫
‫
‫\section{پیاده‌سازی و آزمون}
‫
‫
‫\subsection{پرونده‌های پروژه}
‫
‫پرونده‌های استفاده‌شده در این آزمایش عبارت‌اند از:
‫
‫\شروع{فقرات}
‫\فقره \کد{comparator\_1bit.v}: پیمانه‌ی یک‌بیتی با سه انتساب پیوسته؛
‫\فقره \کد{comparator\_4bit.v}: پیمانه‌ی چهاربیتی شامل چهار نمونه‌ی یک‌بیتی؛
‫\فقره \کد{serial\_comparator.v}: مقایسه‌کننده‌ی دنباله‌ای تک‌پیمانه‌ای؛
‫\فقره \کد{tb\_comparator\_4bit.v}: آزمون جامع مقایسه‌کننده‌ی چهاربیتی؛
‫\فقره \کد{tb\_serial\_comparator.v}: آزمون جامع مقایسه‌کننده‌ی دنباله‌ای.
‫\پایان{فقرات}
‫
‫ساختارهای رفتاری \کد{initial}، حلقه و \کد{task} فقط در آزمون‌گرها استفاده شده‌اند و جزئی
‫از سخت‌افزار طراحی‌شده نیستند. سه پرونده‌ی اصلی مدار شامل هیچ \کد{always}، \کد{initial}،
‫\کد{case}، تابع یا وظیفه‌ای نیستند.
‫
‫
‫\subsection{محیط اجرا}
‫
‫شبیه‌سازی با نسخه‌ی ۱۲٫۰ نرم‌افزار \لر{Icarus Verilog} نصب‌شده در مسیر سیستمی
‫\کد{/usr/bin} انجام شد. گزینه‌ی \کد{-Wall} برای نمایش هشدارها و گزینه‌ی \کد{-g2012}
‫برای انتخاب زبان وریلاگ/سیستم‌وریلاگ نسخه‌ی ۲۰۱۲ به کار رفت.
‫
‫
‫\subsection{آزمون مقایسه‌کننده‌ی چهاربیتی}
‫
‫دو حلقه‌ی تو در تو تمام مقادیر صفر تا ۱۵ را برای هر یک از ورودی‌های $a$ و $b$ تولید
‫می‌کنند. بنابراین تعداد کل مقایسه‌ها برابر است با:
‫
‫\شروع{equation}
‫16 \times 16 = 256
‫\پایان{equation}
‫
‫پس از اعمال هر زوج ورودی، خروجی‌های مدار با نتایج عملگرهای مرجع \کد{>}، \کد{==} و
‫\کد{<} مقایسه می‌شوند. در صورت مشاهده‌ی اختلاف، آزمون‌گر ورودی و خروجی اشتباه را چاپ و
‫شبیه‌سازی را متوقف می‌کند. در اجرای انجام‌شده هیچ اختلافی مشاهده نشد.
‫
‫\شروع{شکل}[ht]
‫\تنظیم‌ازوسط
‫\centerimg{3.png}{16cm}
‫\شرح{کامپایل و عبور موفق تمام ۲۵۶ حالت مقایسه‌کننده‌ی چهاربیتی}
‫\برچسب{شکل:نتیجه آزمون چهاربیتی}
‫\پایان{شکل}
‫
‫
‫\subsection{آزمون مقایسه‌کننده‌ی دنباله‌ای}
‫
‫در آزمون دنباله‌ای نیز تمام ۲۵۶ زوج عدد چهاربیتی بررسی می‌شوند. پیش از ارسال هر زوج، وظیفه‌ی
‫\کد{do\_reset} حالت مدار و خروجی \کد{010} را کنترل می‌کند. سپس حلقه‌ی داخلی از بیت شماره‌ی
‫۳ تا بیت شماره‌ی صفر حرکت می‌کند. پس از هر لبه‌ی بالارونده‌ی ساعت، خروجی واقعی با نتیجه‌ی
‫پیشوندی که تا همان لحظه دریافت شده است مقایسه می‌شود. بنابراین علاوه بر ۲۵۶ کنترل بازنشانی،
‫تعداد بررسی‌های پیشوندی برابر است با:
‫
‫\شروع{equation}
‫256 \times 4 = 1024
‫\پایان{equation}
‫
‫این روش حالت برابری، تفاوت در هر یک از چهار موقعیت و حفظ نتیجه پس از مشاهده‌ی نخستین بیت
‫متفاوت را پوشش می‌دهد. تمام بررسی‌ها با موفقیت عبور کردند.
‫
‫\شروع{شکل}[H]
‫\تنظیم‌ازوسط
‫\centerimg{4.png}{16cm}
‫\شرح{کامپایل و عبور موفق تمام حالت‌ها و پیشوندهای مقایسه‌کننده‌ی دنباله‌ای}
‫\برچسب{شکل:نتیجه آزمون دنباله‌ای}
‫\پایان{شکل}
‫
‫
‫\subsection{تصویر شکل موج}
‫
‫برای نمایش دیداری رفتار ترتیبی، یک اجرای کوتاه از مثال $1010_2$ و
‫$1001_2$ در نمایش‌دهنده‌ی شکل موج انجام شود. تصویر باید سیگنال‌های \کد{clk}، \کد{reset}،
‫$a$، $b$، \کد{gt}، \کد{eq} و \کد{lt} را نشان دهد. در ابتدای تصویر، \کد{eq} باید پس از
‫بازنشانی یک باشد. پس از دو جفت بیت برابر نیز مقدار آن باید حفظ شود. در نبض سوم \کد{gt}
‫یک می‌شود و با وجود متفاوت بودن بیت چهارم، مقدار بزرگ‌تر بودن باید ثابت بماند.
‫
‫\شروع{شکل}[ht]
‫\تنظیم‌ازوسط
‫\centerimg{5.png}{15cm}
‫\شرح{شکل موج بازنشانی و مقایسه‌ی دنباله‌ای مرحله‌به‌مرحله}
‫\برچسب{شکل:شکل موج دنباله‌ای}
‫\پایان{شکل}
‫
‫
‫\section{دستورهای اجرا}
‫
‫ابتدا در پوشه‌ی پروژه، پوشه‌ی خروجی ساخته می‌شود:
‫
\begin{dsdlcodeblock}
\begin{verbatim}
mkdir -p .build
\end{verbatim}
\end{dsdlcodeblock}
‫
‫دستورهای کامپایل و اجرای آزمون مقایسه‌کننده‌ی چهاربیتی عبارت‌اند از:
‫
\begin{dsdlcodeblock}
\begin{verbatim}
iverilog -g2012 -Wall -s tb_comparator_4bit \
  -o .build/comparator4.out \
  comparator_1bit.v comparator_4bit.v tb_comparator_4bit.v

vvp .build/comparator4.out
\end{verbatim}
\end{dsdlcodeblock}
‫
‫خروجی مورد انتظار برابر است با:
‫
\begin{dsdlcodeblock}
\begin{verbatim}
PASS: all 256 four-bit comparisons
\end{verbatim}
\end{dsdlcodeblock}
‫
‫دستورهای کامپایل و اجرای آزمون مقایسه‌کننده‌ی دنباله‌ای عبارت‌اند از:
‫
\begin{dsdlcodeblock}
\begin{verbatim}
iverilog -g2012 -Wall -s tb_serial_comparator \
  -o .build/serial.out \
  serial_comparator.v tb_serial_comparator.v

vvp .build/serial.out
\end{verbatim}
\end{dsdlcodeblock}
‫
‫خروجی مورد انتظار برابر است با:
‫
\begin{dsdlcodeblock}
\begin{verbatim}
PASS: all 256 serial comparisons at every prefix
\end{verbatim}
\end{dsdlcodeblock}
‫
‫
‫\section{نتایج آزمایش}
‫
‫نتایج نهایی در جدول~\رجوع{جدول:نتایج نهایی} خلاصه شده‌اند.
‫
‫
‫\شروع{لوح}[ht]
‫\تنظیم‌ازوسط
‫\شرح{خلاصه‌ی نتایج شبیه‌سازی}
‫
‫\شروع{جدول}{|c|c|c|c|}
‫\خط‌پر
‫\سیاه مدار & \سیاه نوع طراحی & \سیاه تعداد بررسی & \سیاه نتیجه \\
‫\خط‌پر \خط‌پر
‫مقایسه‌کننده‌ی چهاربیتی & ترکیبی و سلسله‌مراتبی & ۲۵۶ زوج ورودی & موفق \\
‫مقایسه‌کننده‌ی دنباله‌ای & ترتیبی و تک‌پیمانه‌ای & ۱۰۲۴ پیشوند و ۲۵۶ بازنشانی & موفق \\
‫\خط‌پر
‫\پایان{جدول}
‫
‫\برچسب{جدول:نتایج نهایی}
‫\پایان{لوح}
‫
‫کامپایل هر دو طراحی با گزینه‌ی نمایش هشدارها انجام شد و پس از افزودن مقیاس زمانی یکسان
‫به پیمانه‌ها، هیچ هشدار کامپایل باقی نماند. آزمون جامع نشان داد که خروجی‌های \کد{gt}،
‫\کد{eq} و \کد{lt} در تمام حالت‌ها یک‌داغ\پاورقی{One-hot} هستند؛ یعنی در هر لحظه دقیقاً
‫یکی از آن‌ها مقدار یک دارد.
‫
‫
‫\section{جمع‌بندی}
‫
‫در بخش نخست آزمایش، معادلات ساده‌ی یک مقایسه‌کننده‌ی یک‌بیتی به کمک انتساب پیوسته پیاده‌سازی
‫شدند. سپس چهار نمونه از این پیمانه با رعایت اولویت بیت پرارزش‌تر به یکدیگر متصل شدند و یک
‫مقایسه‌کننده‌ی چهاربیتی کاملاً ترکیبی و سلسله‌مراتبی تشکیل دادند. در بخش دوم، نتیجه‌ی مقایسه
‫در یک حالت دوبیتی ذخیره شد تا مدار بتواند بیت‌ها را به صورت دنباله‌ای دریافت کند. حالت اولیه
‫برابری است و نخستین جفت بیت متفاوت از سمت پرارزش‌تر نتیجه‌ی نهایی را تعیین می‌کند.
‫
‫محدودیت استفاده‌ی انحصاری از \کد{assign} باعث شد ذخیره‌سازی دنباله‌ای با بازخورد پیوسته‌ی
‫ارباب--برده مدل شود. این ساختار در شبیه‌ساز هدف آزمایش به درستی کار کرد، هرچند برای یک طراحی
‫قابل سنتز صنعتی، استفاده از توصیف ترتیبی لبه‌ای روش متعارف‌تری است. در نهایت، آزمون جامع
‫تمام ورودی‌های ممکن صحت هر دو مدار و نتیجه‌ی هر مرحله از مقایسه‌ی دنباله‌ای را تأیید کرد.
‫
